\documentclass[11pt]{article}

\usepackage[margin=1in]{geometry}
\usepackage[T1]{fontenc}
\usepackage{lmodern}
\usepackage{amsmath}
\setlength{\emergencystretch}{2em}

\title{Review Notes for CoEvoP\&R}
\author{}
\date{}

\begin{document}

\maketitle

\noindent\textbf{Organization.} Each manuscript section is reviewed in its own
part. Comments under one part apply only to that manuscript section unless a
cross-section consistency issue is stated explicitly.

\tableofcontents
\clearpage

\part{Abstract Review}

\section{Overall Assessment}

The abstract has a strong overall structure: it moves from the problem, to the
limitations of prior work, to the proposed method, its safety and evaluation
mechanisms, quantitative results, and the final deployable artifact. The
central novelty---searching the differentiable placement objective itself---is
clear and compelling. The length and overall progression are also appropriate.
There are no serious grammatical problems, but the abstract would benefit from
several stylistic clarifications and a few substantive checks before it is
finalized.

\section{Wording and Clarity Comments}

\begin{enumerate}
  \item ``Fixed expert terms'' would read more naturally as ``fixed,
  expert-designed terms.''

  \item ``Scheduled routed promotions condition later proposals'' is accurate
  but difficult to understand without first reading the Method section. A
  plainer formulation, such as ``results from scheduled routing evaluations
  inform subsequent proposals,'' would be clearer.

  \item ``Pass a restricted objective interface'' is awkward. Candidates pass
  validation checks imposed by the restricted interface; they do not literally
  pass the interface itself.

  \item ``Gives cheap timing proxies selection influence'' is grammatically
  possible but dense. It would be clearer to say that the audit ``allows
  inexpensive timing proxies to influence selection.''

  \item The evaluation list is not parallel: ``ChiPBench Nangate45, ICCAD 2015
  Superblue, and ASAP7 cross-node transfer'' mixes benchmarks, design families,
  technology nodes, and an experiment type. A clearer construction is:
  ``We evaluate CoEvoP\&R on the ChiPBench Nangate45 and ICCAD 2015 Superblue
  benchmarks and study cross-node transfer to ASAP7.'' This wording should be
  checked against the exact experimental setup.

  \item ``Plugs directly into DREAMPlace'' may imply that no integration layer
  is required, whereas the Method describes restricted lowering into an
  \texttt{ObjectiveSpec}. ``Executes within DREAMPlace through the restricted
  objective interface'' may be more precise.

  \item The abstract contains several specialized phrases in quick succession,
  including ``scheduled routed promotions,'' ``term-compatibility rules,'' and
  ``perturbation-based proxy audit.'' The proxy audit is worth retaining if it
  is a principal contribution, but these mechanisms should be described as
  intuitively as possible in the abstract.
\end{enumerate}

\section{Substantive Concerns}

\begin{enumerate}
  \item \textbf{The reported evidence does not fully match the opening
  motivation.} The abstract motivates the work using routed wirelength,
  routing congestion, and post-route timing, but reports improvements only in
  HPWL, density overflow, and placement-RC timing. A reviewer may immediately
  ask whether routing feedback improves routed outcomes. If a strong post-GRT
  result is available, at least one routed-quality result should be reported in
  the abstract.

  \item \textbf{The scope of the numerical claims is ambiguous.} It is unclear
  whether the reported 16.9\%, 36.7\%, 0.70~ns, and 912~ns improvements are
  averages over designs, seeds, or both. In the current sentence, ``on average''
  appears to modify only HPWL and density overflow. The scope of every number
  should be made explicit.

  \item \textbf{The direction of the TNS result should be explicit.} If TNS is
  stored as a negative quantity, ``improving TNS by 912~ns'' can be difficult
  to interpret. Depending on the metric convention, ``reducing TNS magnitude
  by 912~ns'' or another directionally explicit description may be better.

  \item \textbf{The opening correlation claim may be too broad.} The phrase
  ``correlate poorly with routed wirelength, routing congestion, and timing''
  appears to make one general claim about all three metric categories. Unless
  the cited evidence demonstrates weak correlation for every listed metric, a
  more defensible formulation would be ``do not consistently predict'' or
  ``can correlate poorly with.''

  \item \textbf{The benchmark terminology must be exact.} The roles of
  ChiPBench, Nangate45, Superblue, and ASAP7 should be described consistently
  as benchmark suites, design families, technology platforms, or transfer
  settings, as appropriate.

  \item \textbf{The central novelty should remain dominant.} The strongest
  distinction is that prior work tunes parameters, adds fixed terms, or uses
  learned surrogates, whereas CoEvoP\&R searches the complete differentiable
  objective program. This should remain the clearest and most memorable claim
  in the abstract.
\end{enumerate}

The benchmark wording, numerical scope, and routed-quality claim should be
verified against the Experiments and Results section before the abstract is
finalized. These points have not yet been validated because the present review
stops at the end of the Method section.

\section{Proposed Revision}

The following revision improves the wording while preserving the technical
claims currently made in the manuscript. The numerical and benchmark language
should still be checked against the experiments, and a routed-quality result
should be added if available.

\begin{quote}
Analytical placers typically optimize half-perimeter wirelength (HPWL) and
density, although recent post-route benchmarks show that these proxies can
correlate poorly with routed wirelength, routing congestion, and timing.
Existing routing-aware placers address this mismatch using fixed,
expert-designed terms or learned black-box surrogates, while leaving the form
of the differentiable objective unchanged. CoEvoP\&R instead searches over
complete replacement objectives. A large language model (LLM) proposes
readable symbolic objective programs, DREAMPlace evaluates their behavior
within the optimization loop, and results from scheduled routing evaluations
inform subsequent proposals. Each candidate passes a restricted objective
interface, term-compatibility checks, autograd validation, and a cost-scaled
placement, timing, and routing evaluation cascade. A perturbation-based audit
allows inexpensive timing proxies to influence selection only when they
provide information beyond HPWL and agree with downstream timing. We evaluate
CoEvoP\&R on the ChiPBench Nangate45 and ICCAD 2015 Superblue benchmarks and
study cross-node transfer to ASAP7. On the primary ChiPBench benchmark,
CoEvoP\&R reduces HPWL by 16.9\% and density overflow by 36.7\% on average
relative to native DREAMPlace, while improving OpenTimer placement-RC WNS by
0.70~ns and TNS by 912~ns. Each search produces a compact symbolic objective
that can be inspected, edited, and executed within DREAMPlace through the
restricted objective interface.
\end{quote}

\clearpage
\part{Introduction Review}

\section{Review Scope}

The perspectives below are lenses used to evaluate the Introduction; they are
not all topics that must be added to the Introduction itself. The Introduction
should primarily establish the problem, summarize the relevant limitations of
prior work, define the research gap, introduce CoEvoP\&R, state its
contributions, and summarize its most important evidence. Detailed safety,
reproducibility, archive, and evaluation procedures belong in the Method and
Experiments sections.

This review intentionally ignores the keyword list and reference formatting.

\section{Overall Assessment}

The Introduction has a sound narrative arc: proxy mismatch, limitations of
existing routing-aware approaches, objective-program search, contributions,
and results. The main idea is understandable without first reading the Method
section. It does not require restructuring from scratch, but several claims
should be made more precise and the reported evidence should better match the
routing-centered motivation.

\section{EDA and Technical-Accuracy Perspective}

\begin{enumerate}
  \item Calling both HPWL and density ``proxies'' is debatable. HPWL is a
  wirelength proxy, whereas density is primarily a feasibility or overlap
  control objective. A more precise opening would be:

  \begin{quote}
  Modern analytical placers make global placement tractable by minimizing a
  smooth wirelength surrogate together with a density objective.
  \end{quote}

  \item The statement that ChiPBench shows poor correlation with routed
  wirelength, congestion, WNS, TNS, and DRC count should be checked carefully
  against what that work demonstrates for each metric. If the evidence is not
  uniform across all metrics, ``does not consistently predict'' or ``can
  correlate poorly with'' is safer than a categorical claim.

  \item ``The analytical placer optimizes one target while the completed flow
  depends on another'' is memorable but overly singular. The completed flow
  depends on several downstream objectives. A more precise alternative is:

  \begin{quote}
  Consequently, the placer optimizes objectives that do not fully represent
  the quality metrics of the completed physical-design flow.
  \end{quote}

  \item The two-part classification of routing-aware placement is useful, but
  it should not imply that every method fits perfectly into either a
  hand-designed symbolic term or a learned surrogate. The wording should
  present these as two important lines of work rather than an exhaustive
  taxonomy.

  \item ``Require expert redesign for each technology or benchmark family'' is
  too absolute. Hand-designed objectives can transfer, although they may
  require expert retuning or redesign. The latter formulation is easier to
  defend.

  \item ``Black-box models that engineers cannot inspect or edit'' is also too
  strong. Engineers can technically inspect architectures and modify learned
  models. The intended contrast is interpretability and direct editability:

  \begin{quote}
  Learned surrogates provide downstream awareness but are generally less
  interpretable and directly editable than compact symbolic objectives.
  \end{quote}
\end{enumerate}

\section{LLM and Program-Search Perspective}

\begin{enumerate}
  \item The central novelty should be stated as precisely as possible:
  CoEvoP\&R searches the functional form of the differentiable objective while
  keeping the placement optimizer fixed and evaluating each candidate within
  the optimizer's gradient path. This is stronger and clearer than merely
  saying that an LLM is used for placement.

  \item ``Complete symbolic replacement objectives'' may sound like
  unrestricted program synthesis. The Method later shows that programs use a
  restricted grammar, an allowed observable library, and required wirelength
  and density anchors. The Introduction should establish this boundary early,
  for example:

  \begin{quote}
  CoEvoP\&R searches compositions and nonlinear interactions among a
  restricted set of differentiable placement observables.
  \end{quote}

  \item ``This embedding makes optimizer compatibility part of the search
  space'' is imprecise. Compatibility is a validity requirement or evaluation
  constraint, not necessarily a dimension of the search space. A better
  formulation is:

  \begin{quote}
  Embedding candidates in the optimization loop makes optimizer compatibility
  an explicit validity requirement.
  \end{quote}

  \item The Introduction should explain briefly why an LLM is suitable for
  this task: the search space contains readable program structures and
  nonlinear compositions, while automatic EDA measurements provide concrete
  revision feedback.

  \item The term ``co-evolving'' should be explained. At present, the objective
  programs evolve while routing measurements provide feedback; routing itself
  does not appear to form a second evolving population. One sentence should
  clarify what the name CoEvoP\&R is intended to denote.
\end{enumerate}

\section{Conference-Reviewer and Adversarial Perspective}

\begin{enumerate}
  \item The distinction from prior LLM-based placement work must remain
  explicit. Prior work searches parameters, optimizer components, or placement
  actions; CoEvoP\&R searches the differentiable objective program itself.

  \item The Introduction should avoid winning the novelty argument by
  overstating weaknesses of prior work. The phrases most likely to attract
  reviewer objections are ``cannot inspect or edit,'' ``require expert
  redesign,'' and any implication that all previous differentiable objectives
  are entirely fixed or unsearchable.

  \item The claim that the method improves ``post-GRT quality'' needs direct
  support in the Introduction. The reported summary currently includes HPWL,
  density overflow, and placement-RC timing, but no routed metric. Since
  routing feedback appears in both the motivation and the paper title, at least
  one post-GRT wirelength, congestion, overflow, or timing result should be
  reported if available.

  \item The Introduction should distinguish placement-RC timing from
  post-route timing. A reviewer may object if the former is used as evidence
  for a broad post-route timing claim.

  \item The contribution list must contain claims that are individually
  verifiable later in the paper. In particular, ``discovers readable formulas
  that improve placement and post-GRT quality'' should correspond to explicit
  downstream experiments.
\end{enumerate}

\section{Experimental-Methodology Perspective}

\begin{enumerate}
  \item The final numerical sentence should state what is averaged: designs,
  seeds, or both. It should also clarify whether ``on average'' applies to the
  WNS and TNS improvements as well as HPWL and density overflow.

  \item The direction of the 912~ns TNS result should be explicit. If TNS is
  represented as a negative quantity, consider reporting a reduction in TNS
  magnitude or otherwise stating the sign convention.

  \item The baseline wording should be precise. ``Relative to native
  DREAMPlace'' is clearer than ``reduces DREAMPlace HPWL,'' because HPWL is the
  metric and DREAMPlace is the baseline tool or objective.

  \item The result summary should reflect the principal scientific claim. If
  routing-aware objective evolution is the key contribution, a routed outcome
  should accompany the placement metrics.

  \item The Introduction need not describe the full evaluation cascade, but it
  should avoid language that implies final-panel information influences the
  search if the Method explicitly keeps search and final evidence separate.
\end{enumerate}

\section{General-Reader and Language Perspective}

\begin{enumerate}
  \item The second paragraph contains three individual prior-work result claims
  for LaMPlace, GOALPlace, and RoutePlacer. These establish that prior methods
  are effective, but they interrupt the Introduction's argument and may fit
  better in Related Work. A compact alternative is:

  \begin{quote}
  Learned approaches such as LaMPlace, GOALPlace, and RoutePlacer have
  demonstrated substantial timing and routability improvements.
  \end{quote}

  \item ``Later congestion and density controls build on the same idea'' is
  vague. It should either identify the relevant development more concretely or
  be removed.

  \item The third paragraph introduces too many mechanisms in rapid succession:
  program evolution, gradient-path execution, static admission, autograd
  checks, runtime guards, scheduled routing promotion, and proxy auditing.
  Preserve the objective-search distinction and defer lower-level mechanisms
  to the Method section.

  \item ``Scheduled promotions return routed evidence to later proposals'' and
  ``gives cheap timing signals selection influence'' are difficult for a
  general reader. A clearer summary is:

  \begin{quote}
  Placement, timing, and scheduled routing measurements are returned as
  feedback for subsequent proposals.
  \end{quote}

  \item The first and fourth contributions overlap because both emphasize
  readable symbolic objectives. The first should focus on the search
  formulation; the fourth should focus on the inspectable, editable artifact
  returned by the search.
\end{enumerate}

\section{Recommended Revision Priorities}

The Introduction is conceptually solid. Revision should prioritize the
following changes:

\begin{enumerate}
  \item Make the characterization of prior work more measured and defensible.
  \item Define the restricted objective-program search space accurately.
  \item State the novelty as search over the objective's functional form with
  the optimizer held fixed.
  \item Explain the meaning of ``co-evolving.''
  \item Add routed experimental evidence that supports the motivation and
  title, if such evidence is available.
  \item Clarify the averaging scope and direction of all reported results.
  \item Reduce implementation detail in the introductory method summary.
\end{enumerate}

\clearpage
\part{Related Work Review}

\section{Review Scope}

This part reviews the substance and organization of the Related Work section.
It focuses on technical accuracy, fairness to prior work, coverage of the most
relevant research directions, the clarity of the novelty comparison, and
readability. Citation and reference formatting are intentionally excluded.

\section{Overall Assessment}

The Related Work section is concise and organized around three useful themes:
routing-aware placement, automated placement search, and executable LLM
program evolution. It includes several of the closest comparisons, especially
RoutePlacer, EvoPlace, VeoPlace, Eureka, and LLM-SR. Each paragraph also tries
to identify the artifact modified by prior work and contrast it with
CoEvoP\&R's objective-program search.

The main weakness is that several technically different methods are compressed
into categories that are too broad. Reviewers familiar with these papers may
challenge the resulting descriptions. The section should make its comparison
axes more precise and devote more attention to the closest work instead of
listing many broad LLM-for-EDA examples.

\section{Technical-Accuracy Perspective}

\begin{enumerate}
  \item \textbf{The historical attribution to RePlAce should be corrected.}
  ``RePlAce established the electrostatic density formulation'' is inaccurate
  or at least misleading. The electrostatics-based placement formulation is
  generally associated with ePlace; RePlAce subsequently advanced this
  analytical-placement framework. A safer formulation is:

  \begin{quote}
  ePlace introduced the electrostatics-based placement formulation, which
  RePlAce and DREAMPlace subsequently advanced through improved optimization
  and GPU acceleration.
  \end{quote}

  If ePlace is not discussed elsewhere, it should likely be added here.

  \item \textbf{RUDY should be described as a routing-demand estimator.} RUDY
  is not inherently a differentiable objective term. It estimates routing
  demand from placement geometry, and later placement methods can incorporate
  RUDY-derived quantities into congestion-aware objectives. A more precise
  formulation is:

  \begin{quote}
  Hand-designed routing estimators such as RUDY approximate routing demand
  from placement geometry and can be incorporated into congestion-aware
  placement objectives.
  \end{quote}

  \item \textbf{LaMPlace, GOALPlace, and RoutePlacer should not be presented as
  though they use equivalent mechanisms.} Their roles should be mapped
  explicitly:
  \begin{itemize}
    \item LaMPlace learns pixel-level masks for cross-stage metrics in macro
    placement.
    \item GOALPlace learns post-route-informed density targets using an
    empirical-Bayes approach.
    \item RoutePlacer trains a graph neural network as a differentiable
    approximation of routability and integrates it into gradient-based
    placement.
  \end{itemize}

  The current phrase ``through masks, graph neural networks, and learned
  routability signals'' does not clearly associate each mechanism with its
  method.

  \item \textbf{The comparison with LaMPlace should acknowledge task scope.}
  LaMPlace concerns macro placement and generates masks for sequential
  placement decisions, whereas CoEvoP\&R operates inside a mixed-size
  analytical placement objective. Treating both simply as learned downstream
  surrogates hides this important distinction.

  \item \textbf{The first heading does not match all of its contents.}
  ``Routability-Aware Objectives'' begins with RePlAce and DREAMPlace, which
  are foundational analytical placers rather than specifically
  routability-aware objective methods. A clearer heading is
  ``Analytical and Routability-Aware Placement.''

  \item \textbf{The second heading is too narrow.} ``Foundation-Model
  Placement Search'' also includes reinforcement learning and Bayesian
  optimization, neither of which is necessarily foundation-model search. A
  better heading is ``Automated and Foundation-Model Placement Search.''
\end{enumerate}

\section{Novelty-Positioning Perspective}

\begin{enumerate}
  \item \textbf{EvoPlace deserves a sharper comparison.} It appears to be one
  of the closest works in terms of LLM-guided evolution around analytical
  placement. The essential distinction should be stated directly:

  \begin{quote}
  EvoPlace searches optimizer implementation components while retaining the
  placement objective; CoEvoP\&R fixes the optimizer and searches the objective
  program.
  \end{quote}

  \item \textbf{Eureka is also methodologically close.} Its generated reward
  functions execute during a learning loop, so the statement ``These systems
  evaluate generated programs after execution'' is too broad. The distinction
  is not simply execution after the fact versus execution inside an optimizer.
  A more defensible comparison is:

  \begin{quote}
  Unlike prior program-evolution systems, CoEvoP\&R must admit generated
  objective expressions into an existing differentiable EDA optimization loop,
  requiring finite values, valid gradients, and stable behavior throughout
  placement.
  \end{quote}

  \item The novelty should be organized around the \emph{artifact modified}:
  parameters, macro-placement actions, optimizer components, learned
  surrogates, or the differentiable objective form. This is the cleanest basis
  for distinguishing CoEvoP\&R from the closest methods.

  \item The section should avoid implying that using an LLM or executing a
  generated program is itself the central novelty. The stronger claim is the
  search over restricted, differentiable objective programs evaluated inside
  the fixed analytical optimizer and informed by downstream routing evidence.
\end{enumerate}

\section{Consistency with the Method Section}

\begin{enumerate}
  \item ``CoEvoP\&R places generated code inside DREAMPlace's autograd path''
  conflicts slightly with the Method statement that raw model code is never
  imported or executed. The Related Work section should refer to ``lowered
  objective expressions'' or ``validated symbolic programs.''

  \item ``Before routed validation'' suggests that every valid objective is
  routed. According to the Method, only scheduled Tier-C promotions receive
  routed evaluation. A more accurate formulation is:

  \begin{quote}
  before placement evaluation and possible promotion to routed validation.
  \end{quote}

  \item ``Scheduled routed promotions condition later objective proposals'' is
  unnecessarily method-specific and difficult to parse. ``Uses scheduled
  routing results as feedback for subsequent objective proposals'' is clearer.

  \item Terminology should remain consistent across sections, particularly
  ``routing-aware'' versus ``routability-aware,'' ``post-GRT'' versus
  ``post-route,'' and ``generated code'' versus ``lowered objective program.''
\end{enumerate}

\section{Coverage and Balance Perspective}

\begin{enumerate}
  \item The routing-aware paragraph moves directly from RUDY to three learned
  methods. This may appear selective to a domain reviewer. It should briefly
  acknowledge other important classical directions, such as cell inflation,
  congestion-aware density control, global-routing feedback, pin-density or
  pin-accessibility terms, and differentiable congestion estimation.

  \item These classical approaches do not all require detailed discussion. A
  single synthesis sentence can demonstrate awareness of the broader field and
  prevent the two-category framing from appearing constructed only to motivate
  the proposed method.

  \item The broad LLM-for-EDA examples---ChipNeMo, RTLCoder, MasterRTL, and
  VeriLoC---are less directly relevant than EvoPlace, Eureka, LLM-SR, or
  RoutePlacer. If space is limited, shorten the general survey and use the
  recovered space for deeper comparison with the closest approaches.

  \item ChiPBench is an evaluation benchmark rather than a placement-search
  method. Its role in validating whether placement gains survive downstream
  flow is relevant, but this role should be separated clearly from the list of
  systems that search parameters, optimizer components, or actions.
\end{enumerate}

\section{General-Reader and Language Perspective}

\begin{enumerate}
  \item Several paragraphs read as lists of papers followed by one sentence
  about CoEvoP\&R. More synthesis would help the reader understand the design
  dimensions and tradeoffs instead of merely recognizing method names.

  \item Ending every paragraph with a short CoEvoP\&R contrast is useful but
  repetitive. The paper could synthesize the comparisons in a concluding
  paragraph or compact table, allowing the individual paragraphs to describe
  prior work more neutrally.

  \item ``Readable symbolic objective forms'' is somewhat subjective.
  ``Compact symbolic objective forms'' or ``restricted symbolic objective
  programs'' is more measurable and consistent with the Method.

  \item The mechanism descriptions should map explicitly to individual works
  instead of relying on the order of method names and technical phrases.
\end{enumerate}

\section{Recommended Organization}

A clearer three-part organization would be:

\begin{enumerate}
  \item \textbf{Analytical and routability-aware placement.} Cover
  foundational analytical objectives, hand-designed routing estimators,
  congestion-aware techniques, and learned downstream guidance.

  \item \textbf{Automated search for placement.} Organize methods by the
  artifact searched: parameters, placement actions, optimizer code, or
  objective form.

  \item \textbf{Executable program evolution.} Discuss reward-function,
  heuristic, symbolic-equation, and code evolution, followed by the
  differentiability and EDA-validation requirements specific to CoEvoP\&R.
\end{enumerate}

If space permits, a compact comparison table could make the principal
distinctions easier to verify:

\begin{center}
\small
\begin{tabular}{|p{0.20\textwidth}|p{0.23\textwidth}|p{0.22\textwidth}|p{0.22\textwidth}|}
\hline
\textbf{Method family} & \textbf{Artifact modified} & \textbf{Feedback} &
\textbf{Output} \\
\hline
AutoDMP and\newline HyperPlace & Hyperparameters & Placement metrics & Parameter
configuration \\
\hline
VeoPlace & Macro actions & Placement metrics & Placement \\
\hline
EvoPlace & Optimizer components & HPWL & Optimizer code \\
\hline
RoutePlacer & Learned routability model & Routing labels & Neural surrogate \\
\hline
CoEvoP\&R & Differentiable objective form & Placement, timing, and routing &
Symbolic objective \\
\hline
\end{tabular}
\end{center}

\section{Recommended Revision Priorities}

\begin{enumerate}
  \item Correct the historical attribution of electrostatics-based placement
  and consider adding ePlace.
  \item Describe RUDY as a routing-demand estimator rather than inherently as
  an objective term.
  \item Separate the mechanisms and scopes of LaMPlace, GOALPlace, and
  RoutePlacer.
  \item Deepen the comparisons with EvoPlace and Eureka.
  \item Replace ``generated code'' with terminology consistent with the
  restricted lowering interface.
  \item Broaden the classical routing-aware coverage with one concise synthesis
  sentence.
  \item Reduce broad LLM-for-EDA examples if space is needed for closer work.
\end{enumerate}

\clearpage
\part{Method Review}

\section{Review Scope}

This part reviews the Method section, Algorithm~1, Figure~1, and Figure~2. It
separates fundamental logical inconsistencies from lower-priority notation,
reproducibility, and presentation issues. The review does not use or assess the
Experiments and Results section.

\section{Overall Assessment}

The Method has a strong conceptual architecture. Its most promising elements
are complete objective replacement, a restricted symbolic language,
optimizer-embedded validation, cost-aware evaluation, negative memory, and an
auditable archive. The separation between the LLM, the controller, and the EDA
tools is also valuable.

The principal problem is not a lack of technical depth. Several sophisticated
mechanisms are described in parallel without all of their mathematical and
algorithmic connections being closed. These inconsistencies should be resolved
before sentence-level proofreading because they affect the core claims of the
paper.

\section{Critical Internal Inconsistencies}

\subsection{Figure 1 Contradicts Complete Objective Replacement}

The Method explicitly states that Equation~2 defines the complete minimized
objective and is not a residual around a hidden native loss. Figure~1 instead
depicts

\[
f_\theta(\mathbf{x};d)=f_{\mathrm{base}}(\mathbf{x};d)
+\sum_j\alpha_jh_j(T(\mathbf{x};d)),
\]

which is a residual added to a base objective. This directly contradicts one
of the paper's central novelty claims. Figure~1 should instead visualize the
restricted complete objective, for example,

\[
f_\theta(\mathbf{x};d)=\kappa_{\mathrm{wl}}
F_\theta\!\left(\widetilde{\mathcal T}(\mathbf{x};d)\right).
\]

This is the most urgent figure correction.

\subsection{The Timing-Proxy Weight Does Not Affect Selection}

The Method defines a timing-proxy weight $\gamma_q$ in the proxy-admission
equation, but $\gamma_q$ never appears in the Tier-A selection score, final
score, parent policy, or archive-update equation. Therefore, the mathematical
method does not currently support the claim that timing proxies affect
selection.

The timing contribution must enter an explicitly defined score, eligibility
rule, ranking, or parent policy. For example, if consistent with the intended
implementation, a term of the following form could be incorporated into a
selection score:

\[
\sum_q \gamma_q w_q g_q(f).
\]

The algorithm introduces a second ambiguity by running Tier~B ``only if
proxy-admitted.'' Tier-B evaluation appears necessary to obtain the proxy
values being audited. The paper must distinguish admission of a global proxy
signal from admission of an individual candidate to Tier~B; otherwise, the
logic appears circular.

\subsection{Routed Feedback Does Not Formally Affect Parent Selection}

The text states that each routed promotion can change parent selection,
retrieval, or failure context. However, the parent probability depends only on
the Tier-A score, distance from a fixed mechanism, and mechanism frequency. No
Tier-C or routed quantity appears in the parent policy.

Routed results can clearly enter later prompts through feedback
reconstruction, but their effect on parent selection is not defined. The paper
should choose one of two interpretations:

\begin{enumerate}
  \item State that routed evidence affects retrieval and proposal context but
  not parent selection; or
  \item Add a routed-evidence term, eligibility update, or archive update that
  explicitly changes the parent pool or parent weights.
\end{enumerate}

Because routing feedback is central to the title, this connection must be
unambiguous.

\subsection{The Search and Final Evidence Split Is Ambiguous}

The text says final-panel metrics never enter prompts or archive fitness, but
the returned objective is selected using $S_{\mathrm{final}}$ over the routed
subset of the final archive. It is unclear whether ``final'' means the archive
at generation $G$, search-panel Tier-C evidence, or a held-out final evaluation
panel.

If held-out final-panel metrics select $f^*$, the panel is no longer a purely
held-out reporting panel. A clearer formulation would define separate panels,

\[
\mathcal D_S=\text{search panel},\qquad
\mathcal D_F=\text{held-out final panel},
\]

select using search-panel routed evidence,

\[
f^*=\arg\max_{f\in\mathcal A_G^{C,S}}S_C^S(f),
\]

and evaluate the already selected $f^*$ on $\mathcal D_F$ only for final
reporting.

\subsection{The Complexity Tie-Break Claim Is Mathematically Incorrect}

The final score subtracts $\eta_Cc(f)$ from the routed-rank score and then says
that complexity matters only when ranks tie. This is not generally true: the
complexity term can reverse a non-tied rank ordering unless $\eta_C$ is bounded
appropriately.

If complexity is intended only as a tie-breaker, use a lexicographic rule:

\begin{enumerate}
  \item Maximize the routed-rank score.
  \item Among tied candidates, minimize complexity.
\end{enumerate}

Alternatively, provide a bound proving that the complexity penalty cannot
overcome any nonzero rank difference.

\section{Problem Formulation}

\begin{enumerate}
  \item Clarify whether each term scale $\kappa_j$ is specific to a design,
  seed, objective, or complete search run.

  \item Specify which placement and objective call determine the ``first-call''
  scale and how zero or extremely small values are handled.

  \item Clarify whether $\kappa_{\mathrm{wl}}$ is design-specific and how it is
  computed.

  \item Define exactly what constitutes a wirelength or density ``anchor.'' It
  is unclear whether an anchor must appear additively or may occur only inside
  a nonlinear interaction.

  \item Explain how the panel-specific metric scales $\nu_j$ are chosen and
  demonstrate that their construction does not use held-out final-panel
  information.

  \item State whether budgets count candidates, complete panels, individual
  designs, or design-seed runs. The mathematical text calls them placement,
  timing, and routing ``calls,'' while the algorithm increments once per
  candidate.

  \item The notation $F_\theta$ may suggest trained parameters. If $\theta$
  represents an LLM-proposed program structure and constants rather than
  learned parameters, say so explicitly.

  \item ``Keeping the optimizer fixed'' should specify which elements remain
  unchanged: initialization, learning-rate schedule, density-weight schedule,
  preconditioning, legalization, detailed placement, and term
  implementations. A reviewer will want to know whether objective replacement
  indirectly changes any optimizer-specific schedules.

  \item Clarify whether WNS and TNS are treated as signed negative quantities
  so that $o_j=+1$ is directionally correct.
\end{enumerate}

\section{Archive-Conditioned Proposal}

The proposal mechanism is promising but requires additional definitions for
reproducibility.

\begin{enumerate}
  \item Define $d_{\mathrm{base}}$, the distance metric, and the set of ``fixed
  mechanisms'' against which distance is measured.

  \item Define a mechanism signature and explain how signatures are extracted
  from the lowered program.

  \item Define the bounds that permit a candidate to enter the near-miss pool.

  \item Give the cardinalities and selection criteria for
  $\operatorname{Top}$, $\operatorname{Diverse}$,
  $\operatorname{NearMiss}$, and $\operatorname{RelevantFail}$.

  \item Explain how relevant failures are matched to a selected parent's terms
  or mechanisms.

  \item State the number and coefficient values of deterministic siblings and
  how their coefficient ranges are selected.

  \item Give the maximum number of provider retries and identify the sampling
  temperature, randomization settings, and decoding limits.

  \item ``Prevents repeated sampling'' is too strong. The mechanism-count
  denominator discourages repeated sampling but does not prevent it.

  \item Distinguish static/schema repair retries from autograd rejection,
  placement failure, and completed metric regression. The phrase ``malformed
  or rejected responses'' could imply that every failure type triggers an
  immediate provider retry, while Figure~2 suggests that some failures enter
  negative memory instead.

  \item Clarify whether the claimed $K$ independent samples remain independent
  after provider repair retries and whether deterministic provider behavior is
  controlled through explicit seeds when available.
\end{enumerate}

\section{Safe Objective Lowering}

The separation between raw responses and lowered expressions is a major
strength. However, the searched language and numerical gates require more
formal specification.

\subsection{Define the Allowed Term Library}

The Method names weighted-average and log-sum-exp wirelength, electric and
bell-shaped density, soft-RUDY, long-net pressure, pin density, and
pin-count-weighted pressure, but does not provide their mathematical
definitions. Because the objective program is the central searched artifact,
the paper should provide a table or appendix containing, for every allowed
term:

\begin{itemize}
  \item Mathematical expression
  \item Required inputs
  \item Smoothing constants
  \item Normalization
  \item Gradient behavior
  \item Computational cost
  \item Availability conditions
\end{itemize}

Without these definitions, the reader cannot reproduce or fully assess the
objective search space.

\subsection{Clarify the Restricted Language}

\begin{enumerate}
  \item ``The action space is finite'' may be too strong if numeric constants
  can take arbitrary values. ``Restricted action space'' or ``bounded program
  structure'' is safer unless the literal set is explicitly finite.

  \item Define ``safe division,'' including its epsilon or clamping rule.

  \item State whether AST depth and primitive limits apply per expression or to
  the complete program, including all local assignments.

  \item Define allowed coefficient ranges and explain how they were selected.

  \item Explain how baseline clones are detected: structural equality,
  canonicalized AST equality, or numerical equivalence.

  \item Define exactly how ``one wirelength family'' and ``one density family''
  are counted when terms occur multiple times or inside nonlinear expressions.
\end{enumerate}

\subsection{Strengthen the Dynamic Safety Definition}

The condition

\[
0<\|\nabla_{\mathbf{x}}f(\mathbf{x}_0)\|_2<\infty
\]

may reject a valid objective if the initialization happens to be stationary.
Conversely, a very large but finite gradient passes. Consider explicit lower
and upper numerical thresholds, possibly normalized against the native
objective.

The paper should also clarify whether the dynamic gate is applied independently
for every design and seed. Because the first-call test establishes only local
executability, the notation $\mathcal F_{\mathrm{safe}}$ may overstate the
guarantee. ``Locally executable candidates'' would be more precise unless a
stronger property is established.

Runtime guards should specify whether they test only finiteness or also detect
finite but exploding objectives and gradients.

\section{Cost-Scaled Evaluation}

The three-tier architecture is sensible, but this subsection contains several
conceptually different mechanisms. It would be clearer if divided into:

\begin{enumerate}
  \item Evaluation cascade
  \item Tier-A feasibility and scoring
  \item Timing-proxy audit
  \item Routed promotion and final selection
\end{enumerate}

\subsection{Evaluation Cascade}

\begin{enumerate}
  \item ``OpenROAD/ChiPBench'' makes ChiPBench sound like a routing engine.
  Describe the relationship precisely, such as an OpenROAD flow configured or
  evaluated through ChiPBench, if that matches the implementation.

  \item Clearly distinguish FLUTE/Elmore timing from OpenTimer placement-RC
  timing and state which one is used during search, candidate selection, and
  final reporting.

  \item State whether placement-RC timing is evaluated before or after
  legalization and which parasitics, libraries, clocks, and constraints are
  used.

  \item Clarify whether every promoted Tier-C candidate is evaluated on the same
  designs and seeds so that downstream comparisons are genuinely matched.
\end{enumerate}

\subsection{Tier-A Gate and Score}

\begin{enumerate}
  \item Define ``structural failure'' and ``severe regression,'' including all
  thresholds.

  \item Explain whether worst-design gains first aggregate placement seeds or
  whether the minimum is taken over every design-seed pair.

  \item Define the rank averaged by $\bar r_A$: the ranked quantity, direction,
  treatment of ties, and aggregation across metrics and design-seed cells.

  \item Define the complexity function $c(f)$ and its scale.

  \item ``The strongest HPWL/overflow member'' is ambiguous because a fixed
  portfolio can have a Pareto frontier rather than one strongest member.

  \item Formalize the additional fixed-portfolio and Pareto gate. It currently
  appears only in prose and is not represented in $\chi_A$ or the algorithm.

  \item Explain why HPWL and density enter both as direct gains and through the
  rank term, and whether this intentionally gives them additional weight.
\end{enumerate}

\subsection{Timing-Proxy Audit}

The perturbation audit needs a reproducible definition of the perturbations
and statistics. Specify:

\begin{itemize}
  \item What is perturbed
  \item How perturbations are sampled
  \item Number of samples
  \item Whether correlations are pooled across designs or computed per design
  \item Minimum sample count
  \item Treatment of ties and missing timing values
  \item Exact computation of $a_C(q)$
  \item Meaning of ``preserve rank direction''
  \item Justification for the 0.70 and 0.95 thresholds
\end{itemize}

Low correlation with HPWL establishes only non-redundancy with HPWL; it does
not establish that a timing proxy is informative rather than noisy. The routed
agreement test is therefore essential and should be specified rigorously.

The phrase ``held-out routed promotions'' should also identify what those
promotions are held out from and confirm that they are distinct from the final
reporting panel.

\subsection{Routed Promotion and Final Ranking}

\begin{enumerate}
  \item Define how high-$S_A$ candidates and occupants of uncovered mechanism
  cells are balanced within a promotion batch.

  \item Define ``matched downstream evidence,'' particularly when some
  post-GRT metrics are unavailable.

  \item Define $\bar r_C$: metrics included, orientations, design and seed
  aggregation, treatment of ties, and treatment of missing values.

  \item Ensure that candidates are not rewarded for missing difficult metrics
  merely because the score averages only ``reported'' metrics.

  \item Consider whether a feasibility or worst-case gate is required before
  final rank averaging so that a catastrophic regression in one routed metric
  cannot be hidden by improvements in others.
\end{enumerate}

\section{Archive and Closed-Loop Feedback}

The detailed archive record is a strength and supports the auditability claim.
The MAP-Elites and feedback mechanisms nevertheless require further
specification.

\begin{enumerate}
  \item Define all bin boundaries for objective size, HPWL, density, and timing.

  \item Explain whether mechanism labels are determined from declared term
  usage, parsed AST structure, or observed behavior.

  \item Define the treatment of hybrid mechanisms and missing timing.

  \item Give deterministic tie-breaking rules for candidates with equal scores
  in one cell.

  \item Use island-specific notation in the archive-update equation or explain
  how the global archive and island archives relate.

  \item Define migration frequency, number of migrants, and whether copied
  elites retain their original island identity and lineage.

  \item Avoid overloading the notation $m_j$, $m(f)$, and $\mathbf m_g$ for an
  evaluator metric, mechanism category, and feedback scorecard.

  \item Explain how ``repair lessons'' are generated. State whether they are
  deterministic templates derived from failure codes or model-generated
  summaries.

  \item Clarify which archive fields a routed result can change. This is
  necessary to support the claim that routing evidence changes the closed-loop
  search rather than merely being stored for later display.

  \item ``Provider snapshot'' may be difficult to guarantee for closed model
  APIs. Specify exactly which provider and model-version metadata can be
  archived.
\end{enumerate}

\section{Algorithm 1}

Algorithm~1 is useful, but it does not yet fully match the prose and equations.

\begin{enumerate}
  \item The number of islands $I$ is used but is not listed as an input.

  \item The initial feedback packet $\mathbf m_0$ is not defined.

  \item Dynamic finite-value and autograd admission is not shown.

  \item Provider retries and their budget are omitted.

  \item The lowered candidate set $\mathcal S_g$ defined earlier is never used;
  the algorithm instead introduces raw responses $\mathcal Y_g$ and lowers
  them again.

  \item The Tier-B ``proxy-admitted'' condition must be reconciled with the
  proxy-audit procedure.

  \item The timing weight $\gamma_q$ never affects insertion, scoring, parent
  sampling, or archive updates.

  \item Budget increments should state whether a multi-design panel counts as
  one evaluation or multiple design-seed calls.

  \item Once $B_A$ is exhausted, only the innermost calibration loop breaks;
  the outer candidate and generation loops continue. The intended termination
  behavior should be shown.

  \item Migration scheduling should explicitly use $\tau_{\mathrm{mig}}$.

  \item The algorithm should define behavior when no candidate receives valid
  Tier-C evidence and $\mathcal A_G^C$ is empty.

  \item Held-out final-panel evaluation is absent. Add it after selection if
  the panel is intended only for reporting.

  \item The algorithm should show the fixed-portfolio/Pareto elite gate if that
  gate affects archive insertion or routed eligibility.
\end{enumerate}

Algorithm~1 should be revised only after the timing score, routed-feedback
path, and evidence split are settled; otherwise, it will need to be rewritten
multiple times.

\section{Figure 1 Review}

\subsection{Strengths}

\begin{itemize}
  \item The proposal, tools, and memory stages make the high-level loop
  intuitive.
  \item The closed feedback path is visually prominent.
  \item Negative memory and the MAP-Elites archive are represented clearly.
\end{itemize}

\subsection{Concerns}

\begin{enumerate}
  \item Replace the residual-objective formula, which contradicts complete
  objective replacement.

  \item Tier~C should be labeled ``scheduled routed promotions'' rather than
  ``selected finalist.'' Final selection occurs only after routed evidence is
  collected.

  \item The Tier A $\rightarrow$ Tier B $\rightarrow$ Tier C flow implies that
  every Tier-C candidate passes through Tier~B. This does not clearly match the
  prose, where timing is conditional and Tier-C promotion is based primarily
  on Tier-A elites and mechanism coverage.

  \item Mutation labels such as ``recombine'' and ``simplify'' should appear in
  the text if they are actual controller modes. Otherwise, use only the modes
  implemented and described.

  \item The five-dimensional MAP-Elites archive is drawn as a two-dimensional
  grid. Label it as a schematic projection.

  \item Several labels are likely to be too small at final two-column print
  size. Verify the minimum font size in the compiled paper.

  \item The caption should explain the green, gray, and red paths and identify
  which failures are retried versus stored as negative memory.

  \item ``BRAIN'' may sound anthropomorphic. ``Proposal,'' ``Evaluation,'' and
  ``Archive and Feedback'' would be more technical stage names.

  \item The figure should show explicitly whether routed evidence changes
  parent scoring, archive eligibility, or only the next prompt context.
\end{enumerate}

\section{Figure 2 Review}

\subsection{Strengths}

\begin{itemize}
  \item The figure clearly separates the response schema from the executable
  record.
  \item The rejected-candidate paths are useful.
  \item The example helps make the restricted objective-program concept
  concrete.
\end{itemize}

\subsection{Concerns}

\begin{enumerate}
  \item The caption implies that schema, AST, term, and autograd checks all
  lower a response into \texttt{ObjectiveSpec}. The diagram and Method instead
  place autograd validation after lowering. Distinguish static lowering from
  dynamic admission.

  \item The finite-value and autograd gate occurs on the first DREAMPlace
  objective call, not entirely before DREAMPlace evaluation. The figure should
  show the gate inside or at the boundary of the first DREAMPlace call.

  \item The example term \texttt{soft\_rudy\_pnorm} should match the term names
  and categories used in the prose.

  \item The prompt tuple should use the same field ordering and notation as the
  corresponding equation unless a reason for the difference is stated.

  \item Schema validation appears both in ``Schema-bound JSON'' and in the
  later grammar/schema check. Remove the apparent duplication or explain the
  two distinct schema checks.

  \item Static retry failures and dynamic negative-memory failures should have
  visually distinct destinations.

  \item Text is likely too small when the figure is printed at one-column
  width. Verify readability in the compiled manuscript rather than only in the
  source PDF.

  \item Clarify whether only a lowered \texttt{ObjectiveSpec} enters the
  optimizer, consistent with the statement that raw model code is never
  executed.
\end{enumerate}

\section{Structural and Presentation Recommendations}

\begin{enumerate}
  \item Divide Cost-Scaled Evaluation into explicit subsubsections for the
  cascade, Tier-A score, proxy audit, and routed promotion.

  \item Divide Archive and Closed-Loop Feedback into separate archive and
  feedback subsections.

  \item Avoid placing two section labels on one subsection; give the proxy audit
  and search procedure their own visible headings.

  \item Add a compact notation table. The Method introduces many sets, scores,
  scales, counts, gates, and archive descriptors before all are defined.

  \item Add a term-library table containing the objective observables and their
  formulas.

  \item Include one concrete candidate objective example in the text and trace
  it through lowering, Tier-A evaluation, possible routing promotion, and
  archive insertion.

  \item Reduce forward references such as ``score defined below'' where a short
  reordering would make the procedure easier to follow.
\end{enumerate}

\section{Recommended Revision Order}

\begin{enumerate}
  \item Correct Figure~1's objective formula.
  \item Formalize the search-panel and final-panel split.
  \item Connect $\gamma_q$ to an actual selection rule.
  \item Decide whether routed evidence affects parent selection or only prompt
  context.
  \item Correct the final complexity tie-break.
  \item Define the objective-term library and normalization scales.
  \item Align Algorithm~1 with the corrected logic.
  \item Revise Figure~1 and Figure~2 to match the algorithm precisely.
  \item Only then perform sentence-level condensation and proofreading.
\end{enumerate}

The section's underlying idea remains strong. Once the missing mathematical
connections are made explicit, the Method will become substantially more
convincing, reproducible, and easier to review.

\end{document}
